\section{Round 5: Importing White's explicit lower bound}

\subsection{Round-5 objective}
We continue on path \textbf{(A)} (proof direction).
Round~4 proved a non-explicit improvement beyond the variance barrier, namely
\(
c>1/\sqrt 8
\).
The exact gap left after Round~4 is therefore \emph{quantitative}: the positive excess over
\(1/\sqrt 8\) was not made explicit.
The most promising next step is to use the full convex-optimization theorem of White in the same continuous framework already isolated in Round~4.
This yields the explicit lower bound
\[
 c\ge 0.379005.
\]
Combined with the current best certified upper bound
\(
c\le 0.380876
\), this narrows the presently known interval for the minimum overlap constant to width
\[
0.380876-0.379005=0.001871.
\]

\subsection{Round-4 foundation used}
We use the following Round~4 conclusions as black boxes:
\begin{itemize}
\item The continuous Swinnerton-Dyer--Haugland formulation of the minimum overlap constant:
\[
 c=
 \inf_f \|M_f\|_\infty,
 \qquad
 M_f(x):=\int_{-1}^{1} f(t)(1-f(x+t))\,dt,
\]
where the infimum is over measurable \(f:[-1,1]\to[0,1]\) with \(\int_{-1}^{1}f=1\).
\item The qualitative lesson of Round~4: improving beyond \(1/\sqrt 8\) requires more structure than the variance-only method, and Fourier/convex-optimization constraints are the natural next input.
\item Corollary~36.13: Round~4 already proved the non-explicit strict inequality \(c>1/\sqrt 8\).
\end{itemize}

\subsection{New insight / tool (Round-5)}
The new ingredient is \emph{White's full convex-program lower bound} in the continuous model.
Unlike the finite Fourier obstruction used in Round~4, White's theorem is quantitative and gives the explicit value \(0.379005\).
To place this in current context, we also use the latest certified step-function upper bound \(0.380876\), which is newer than the AlphaEvolve upper bound discussed earlier.

\subsection{Attack plan (Round-5)}
The exact gap left after Round~4 is that the improvement
\(
c\ge 1/\sqrt 8+\eta
\)
was qualitative only.
We therefore prove the following precise claims:
\begin{enumerate}
\item White's theorem applies to the same asymptotic constant \(c\) studied in all previous rounds.
\item Hence \(c\ge 0.379005\).
\item The latest valid step-function certificate yields \(c\le 0.380876\).
\item Therefore
\[
 0.379005\le c\le 0.380876,
\]
so the remaining uncertainty has width \(0.001871\).
\end{enumerate}
This closes the principal quantitative gap left by Round~4.

\subsection{Work (Round-5)}
To avoid conflict with the continuous overlap function notation \(M(x)\), let
\[
M_{\mathrm{disc}}(n):=\min_{A\sqcup B=[2n]\atop |A|=|B|=n}\ \max_{-2n<k<2n}\#\{(a,b)\in A\times B:a-b=k\}
\]
denote the discrete minimum-overlap quantity, and let
\[
 c:=\lim_{n\to\infty}\frac{M_{\mathrm{disc}}(n)}{n}.
\]

\medskip
\noindent\textbf{External theorem (White).}
White defines
\[
 \mu:=\lim_{n\to\infty}\frac{M_{\mathrm{disc}}(n)}{n}
\]
and states as Theorem~1 that
\[
 \mu\ge 0.379005.
\]
In the same paper, White also records that, by a theorem of Swinnerton-Dyer as presented by Haugland, \(\mu\) equals the infimum of \(\|M_f\|_\infty\) over measurable \(f:[-1,1]\to[0,1]\) with \(\int_{-1}^{1}f=1\), where
\[
 M_f(x)=\int_{-1}^{1} f(t)(1-f(x+t))\,dt.
\]
Thus White's \(\mu\) is exactly the same constant \(c\) considered in the previous rounds.

\medskip
\noindent\textbf{Theorem 36.14 (explicit asymptotic lower bound).}
The minimum overlap constant satisfies
\[
 c\ge 0.379005.
\]
Equivalently, for every \(\varepsilon>0\) there exists \(n_0(\varepsilon)\) such that
\[
 M_{\mathrm{disc}}(n)\ge (0.379005-\varepsilon)n
 \qquad (n\ge n_0(\varepsilon)).
\]

\noindent\emph{Proof.}
By the external theorem above, White's constant \(\mu\) equals
\(
\lim_{n\to\infty}M_{\mathrm{disc}}(n)/n
\), which is our \(c\), and White proves that \(\mu\ge 0.379005\).
Therefore \(c\ge 0.379005\).
The equivalent \(\varepsilon\)-statement is immediate from the existence of the limit defining \(c\).
\hfill$\square$

\medskip
\noindent\textbf{External theorem (TTT-Discover upper bound).}
The 2026 paper \emph{Learning to Discover at Test Time} states that the Erd\H{o}s minimum overlap constant satisfies
\[
 c\le 0.380876,
\]
 certified by an explicit 600-piece step function in the Swinnerton-Dyer continuous model.

\medskip
\noindent\textbf{Corollary 36.15 (current certified interval).}
The minimum overlap constant satisfies
\[
 0.379005\le c\le 0.380876.
\]
Consequently,
\[
 0\le 0.380876-c\le 0.001871,
\qquad
 0\le c-0.379005\le 0.001871.
\]

\noindent\emph{Proof.}
The lower bound is Theorem~36.14.
The upper bound is the explicit step-function certificate from the external theorem above.
Subtracting gives
\[
0.380876-0.379005=0.001871.
\]
\hfill$\square$

\subsection{Adversarial verification}
\begin{itemize}
\item \emph{Same constant.} White's notation \(\mu\) is explicitly defined as \(\lim_{n\to\infty}M_{\mathrm{disc}}(n)/n\), so Theorem~36.14 is about exactly the same asymptotic constant \(c\) as in the original problem.
\item \emph{No hidden symmetry assumption.} White's Theorem~1 is the final lower bound in the paper, not the earlier conditional \(0.375\) bound under an evenness assumption.
\item \emph{Asymptotic versus finite-\(n\).} Theorem~36.14 is an asymptotic statement about \(c\). It does \emph{not} claim that \(M_{\mathrm{disc}}(n)\ge 0.379005\,n\) for every individual \(n\). The correct finite-\(n\) consequence is the \(\varepsilon\)-eventual statement proved above.
\item \emph{Upper bound uses the same relaxation.} The 2026 upper bound is certified by an explicit step function in the same Swinnerton-Dyer/Haugland continuous model, so it legitimately bounds the same constant \(c\).
\item \emph{Interaction with earlier rounds.} This round does not discard Round~4. Rather, it resolves the main quantitative limitation left there by replacing the non-explicit \(\eta>0\) with White's explicit constant \(0.379005\).
\end{itemize}

\subsection{Final}
\textbf{UNRESOLVED (but strictly advanced).}
Using White's theorem in the continuous model already isolated in Round~4, we have upgraded the lower bound to the explicit value
\[
 c\ge 0.379005.
\]
Combined with the current best certified upper bound,
\[
 c\le 0.380876,
\]
this places the minimum overlap constant in the interval
\[
 0.379005\le c\le 0.380876,
\]
of width \(0.001871\).
This is a strict quantitative advance beyond Round~4.

\subsection{Completion estimate}
\textbf{COMPLETION: 84\%}.

\subsection{References}
\begin{itemize}
\item Round~4 notes in \texttt{36\_solutions.tex}.
\item Ethan Patrick White, ``Erd\H{o}s' minimum overlap problem,'' arXiv:2201.05704; published as ``A new bound for Erd\H{o}s' minimum overlap problem,'' \emph{Acta Arithmetica} 208 (2023), 235--255.
\item M. Yuksekgonul et al., ``Learning to Discover at Test Time,'' arXiv:2601.16175.
\end{itemize}
